\section{Билет 8. Метод Хука--Дживса}

\ticketq{Метод Хука--Дживса.}

\paragraph{Постановка.}
$f(x) \to \min$, $x \in \R^n$. Метод \textbf{прямой}, или \textbf{нулевого
порядка}: использует только значения $f(x_1), f(x_2), \dots$ и не требует
$\nabla f(x)$, $H(x) = \nabla^2 f(x)$. Строится последовательность
$x^0, x^1, \dots$ с $f(x^{k+1}) < f(x^k)$ по возможности.

\textbf{Параметры:} начальная точка $x^0$; шаги по координатам $h_i > 0$,
$i = \overline{1,n}$; точность $\eps > 0$; коэффициент дробления
$\beta \in (0,1)$. \textbf{Базисная точка} $x^b$ — текущая основная точка,
вначале $x^b = x^0$.

Метод состоит из двух чередующихся шагов: исследующего поиска и поиска по
образцу.

\paragraph{Шаг 1. Исследующий поиск.}
Покоординатная проверка из текущей точки $x$; $e_i$ — $i$-й единичный вектор.
Для каждой координаты $i = \overline{1,n}$:
\[
  \begin{aligned}
    f(x + h_i e_i) &< f(x) &&\Rightarrow\ x := x + h_i e_i ; \\
    \text{иначе } f(x - h_i e_i) &< f(x) &&\Rightarrow\ x := x - h_i e_i ; \\
    \text{иначе} &&&\ \text{координата не меняется.}
  \end{aligned}
\]
Сначала пробуем «вперёд», потом «назад». После всех координат получается
точка $x^e$ — результат исследующего поиска. За один проход каждая координата
меняется не более одного раза, и $f$ только невозрастает.

\paragraph{Шаг 2. Поиск по образцу.}
Если $f(x^e) < f(x^b)$, поиск удачен; направление удачного смещения
$x^e - x^b$, и строится ускоренная точка
\[
  x^p = x^e + (x^e - x^b) = 2x^e - x^b .
\]
Из $x^p$ выполняется \emph{ещё один исследующий поиск}, дающий $x^{e2}$, и
сравнение идёт уже по нему:
\[
  f(x^{e2}) < f(x^e) \ \Rightarrow\ x^b := x^e,\ x^e := x^{e2},\
    \text{повторить поиск по образцу;}
\]
\[
  f(x^{e2}) \ge f(x^e) \ \Rightarrow\ x^b := x^e,\
    \text{вернуться к исследующему поиску.}
\]
Смысл: исследующий поиск за цикл выявил удачное направление — есть смысл
пройти по нему ещё раз, не тратя вычислений на пробные шаги. Это дешёвая
замена одномерной минимизации. Повторный исследующий поиск из $x^p$ нужен
затем, чтобы подправить направление: чистая экстраполяция легко проскакивает
мимо дна оврага.

\paragraph{Условие остановки.}
Если исследующий поиск не дал улучшения, шаги дробятся: $h_i := \beta h_i$.
Метод останавливается при $\max_i h_i < \eps$ — поиск с достаточно мелкой
сеткой уже не улучшает значение.

\paragraph{Теорема (о монотонном убывании).}
\textit{Если новая базисная точка принимается только при строгом уменьшении
$f$, то $f(x^b_{k+1}) \le f(x^b_k)$, а при удачной итерации
$f(x^b_{k+1}) < f(x^b_k)$.}

\textsc{Док-во.} Два случая.

\emph{1) Поиск неудачен:} $f(x^e) \ge f(x^b_k)$ $\Rightarrow$ базисная точка
не меняется, $x^b_{k+1} = x^b_k$ $\Rightarrow$ $f(x^b_{k+1}) = f(x^b_k)$
(дробятся лишь шаги $h_i$).

\emph{2) Поиск удачен:} $f(x^e) < f(x^b_k)$, строится $x^p = 2x^e - x^b_k$.
\[
  f(x^p) < f(x^e) \Rightarrow x^b_{k+1} = x^p :\ f(x^p) < f(x^e) < f(x^b_k) ;
\]
\[
  f(x^p) \ge f(x^e) \Rightarrow x^b_{k+1} = x^e :\ f(x^e) < f(x^b_k) .
\]
В обоих подслучаях $f(x^b_{k+1}) < f(x^b_k)$. $\blacksquare$

\paragraph{Теорема (о конечности исследующего поиска).}
\textit{Один исследующий поиск требует не более $2n$ вычислений $f$.}

\textsc{Док-во.} Для каждой координаты проверяются не более двух точек,
$x + h_i e_i$ и $x - h_i e_i$ $\Rightarrow$ не более двух вычислений на
координату; координат $n$ $\Rightarrow$ не более $2n$. Так как $n$ конечно,
поиск завершается за конечное число операций. $\blacksquare$

\paragraph{Алгоритм.}
\begin{csteps}
  \item Задать $x^0$, $h_i > 0$, $\eps > 0$, $\beta \in (0,1)$; $x^b := x^0$.
  \item Выполнить исследующий поиск из $x^b$, получить $x^e$.
  \item $f(x^e) \ge f(x^b)$: положить $h_i := \beta h_i$; если
        $\max_i h_i < \eps$ — конец, иначе к~шагу~2.
  \item $f(x^e) < f(x^b)$: построить $x^p = 2x^e - x^b$ и выполнить
        исследующий поиск из $x^p$, получив $x^{e2}$.
  \item $f(x^{e2}) < f(x^e)$: положить $x^b := x^e$, $x^e := x^{e2}$ и
        повторить шаг~4. Иначе $x^b := x^e$ и к~шагу~2.
\end{csteps}

\paragraph{Сходимость и ограничения.}
Метод не требует дифференцируемости и применим, когда производные неизвестны
или дороги. Глобального минимума для невыпуклой $f$ он не гарантирует —
сходится к локальному минимуму либо к точке, где выбранные координатные шаги
уже не дают улучшения. Прочие ограничения: много вычислений $f$; скорость ниже
градиентных и ньютоновских методов; сильная зависимость от выбора $h_i$ и от
масштаба переменных (при разных масштабах метод работает медленно); крупные
шаги проскакивают хорошую область, мелкие замедляют сходимость.

\seefull{гл.~2, \S19 (численный пример там же)}
